\section{Coupled deformation and fluid transport}
\label{sec:finite-elements}

Spatially nonuniform drainage couples the constitutive response to fluid
transport and mechanical equilibrium. We implement these balances in
MOOSE~\cite{gaston2009moose}. We consider isothermal, quasistatic
deformation of a saturated body, with no chemical reaction or gravity.
The unknown fields are skeleton displacement \(\mathbf u\) and pore
pressure \(p\). All balance equations are integrated over the reference
mixture domain \(\Omega_0\), with
\(\mathbf F=\mathbf I+\operatorname{Grad}\mathbf u\).
The elastic law supplies total Cauchy stress through
\eqref{eq:total-stress-phase-energy}; the corresponding first Piola stress is
\begin{equation}
 \mathbf P=J\mathbf\sigma\mathbf F^{-T}.
 \label{eq:fe-total-first-piola}
\end{equation}
This stress already contains the fluid-pressure contribution. No additional
effective-stress correction is applied in the momentum residual.

\subsection{Fluid mass and the transport closure}
\label{sec:fe-fluid-closure}

The pore volume per reference mixture volume is
\(J-\phi_{s0}\bar J\). Let \(\bar\rho_f\) be fluid intrinsic density
and \(m_f\) be fluid mass per reference mixture volume. We prescribe a
constant fluid bulk modulus \(K_f>0\), with density \(\bar\rho_{f0}\)
at zero pressure:
\begin{align}
 \bar\rho_f(p)&=\bar\rho_{f0}\exp(p/K_f),
 \label{eq:fe-fluid-eos}\\
 m_f(\mathbf F,p)&=\bar\rho_f(p)
       \left(J-\phi_{s0}\bar J(\mathbf F,p)\right).
 \label{eq:fe-reference-fluid-mass}
\end{align}
The mineral volume is the solution of
\eqref{eq:anisotropic-mineral-eos}, not an independent porosity law.
All mass measures in this section are per unit reference mixture volume,
whereas the intrinsic densities and volume fractions of the preceding sections
are current-configuration measures; the two conventions differ by the factor
\(J\) of \eqref{eq:true-mineral-jacobian}.
The fluid equation of state is an additional constitutive assumption for
the transport problem; it is separate from the elastic mineral law.

For the hydraulic comparison we prescribe constant, isotropic permeability
\(k>0\) in the current configuration and constant fluid viscosity
\(\mu_f>0\). The Darcy volume flux relative to the skeleton is
\(\mathbf q\), measured per current area. Its referential mass flux is
\(\mathbf Q_f\):
\begin{align}
 \mathbf q&=-\frac{k}{\mu_f}\mathbf F^{-T}
                 \operatorname{Grad}p,
 \label{eq:fe-spatial-darcy-flux}\\
 \mathbf Q_f&=J\bar\rho_f\mathbf F^{-1}\mathbf q,
 \label{eq:fe-reference-mass-flux}\\
 &=-\frac{J\bar\rho_f k}{\mu_f}
      \mathbf F^{-1}\mathbf F^{-T}\operatorname{Grad}p.
 \label{eq:fe-reference-darcy-law}
\end{align}
Keeping the spatial permeability isotropic separates the elastic
directional response from anisotropy prescribed directly in the hydraulic
law. A material permeability transported with the skeleton would be a
different constitutive choice.

The same mineral-volume response determines both storage and mechanical
coupling. Differentiating \eqref{eq:fe-reference-fluid-mass} gives
\begin{align}
 \left.\pder{m_f}{p}\right|_{\mathbf F}
 &=\bar\rho_f\left[
       \frac{J-\phi_{s0}\bar J}{K_f}
       -\phi_{s0}\left.\pder{\bar J}{p}\right|_{\mathbf F}
                  \right],
 \label{eq:fe-fluid-storage}\\
 \left.\delta m_f\right|_p
 &=\bar\rho_f J\mathbf B:
                  (\delta\mathbf F\mathbf F^{-1}).
 \label{eq:fe-fluid-deformation-coupling}
\end{align}
The second identity follows from
\eqref{eq:biot-pore-volume-variation}. It supplies an independent check
that the fluid-mass residual uses the same pressure-coupling tensor as
the stress law. At the undeformed zero-pressure state, the total storage
per reference fluid density is
\begin{equation}
 \frac{1}{\bar\rho_{f0}}
 \left.\pder{m_f}{p}\right|_{\mathbf F=\mathbf I,p=0}
 =\frac{1-\phi_{s0}}{K_f}+S_s,
 \label{eq:fe-reference-total-storage}
\end{equation}
where \(S_s\) is the solid storage defined in
\eqref{eq:reference-solid-storage-definition}.

\subsection{Weak balances and plane strain}
\label{sec:fe-weak-balances}

Let \(\mathbf v\) and \(w\) be displacement and pressure test
functions that vanish on their respective essential boundaries.
Mechanical body force \(\mathbf b_0\) and fluid source \(s_f\) are
specified per reference mixture volume. With outward reference normal
\(\mathbf N\), reference traction \(\mathbf t_0\), and outward mass
flux \(\bar Q_f\), the weak balances are
\begin{align}
 0={}&\int_{\Omega_0}\operatorname{Grad}\mathbf v:\mathbf P\,\dd V_0
      -\int_{\Omega_0}\mathbf v\cdot\mathbf b_0\,\dd V_0
      -\int_{\Gamma_t}\mathbf v\cdot\mathbf t_0\,\dd A_0,
 \label{eq:fe-momentum-residual}\\
 0={}&\int_{\Omega_0}w\dot m_f\,\dd V_0
      -\int_{\Omega_0}\operatorname{Grad}w\cdot\mathbf Q_f\,\dd V_0
      +\int_{\Gamma_Q}w\bar Q_f\,\dd A_0
      -\int_{\Omega_0}w s_f\,\dd V_0.
 \label{eq:fe-fluid-residual}
\end{align}
Here \(\Gamma_t\) and \(\Gamma_Q\) are the boundaries carrying
prescribed traction and mass flux, respectively; on \(\Gamma_Q\),
\(\mathbf Q_f\cdot\mathbf N=\bar Q_f\).
The overdot denotes a time derivative following a material point of the
skeleton. Backward Euler
approximates the change of the full reference fluid mass:
\begin{equation}
 \dot m_f^{\,n}\simeq
      \frac{m_f(\mathbf F^n,p^n)-m_f(\mathbf F^{n-1},p^{n-1})}
           {t_n-t_{n-1}}.
 \label{eq:fe-conservative-time-step}
\end{equation}
Thus a sealed, source-free body conserves its integrated fluid mass up to
the nonlinear and discrete balance errors.

Plane strain sets the out-of-plane displacement to zero while retaining
the three-dimensional constitutive law. In particular, \(F_{33}=1\),
and the out-of-plane normal stress is generally nonzero. The factors
\(1/3\) in the logarithmic strain split are unchanged. Mineral volume
is solved locally at each quadrature point on the branch satisfying
\eqref{eq:trace-mineral-stability-domain}, with positive phase volumes.
Its deformation and pressure derivatives follow by implicit
differentiation of \eqref{eq:anisotropic-mineral-eos}.
These derivatives enter both residuals; holding the local root fixed
while differentiating would omit mechanical--hydraulic coupling.

Rigid frictionless loading plates impose a common normal displacement
and zero tangential traction. For force-controlled consolidation, the
common normal displacement is unknown and the integrated plate reaction
equals the prescribed resultant force. This boundary condition permits
the normal traction to redistribute along a plate as drainage proceeds.
It is distinct from prescribing the same normal traction at every
point of that boundary.

\subsection{Reference problems and error measures}
\label{sec:fe-reference-problems}

The isotropic reference is Mandel consolidation between frictionless rigid
plates. We evaluate the pressure and displacement series independently,
using the expressions in Walker et al.~\cite[appendix D]{walker2023poroelasticity}.
The full reference rectangle is \([-a,a]\times[-b,b]\), with
\(a=1\) and \(b=0.1\) in the selected length unit. Its vertical faces are
drained and traction-free; the plates are impermeable. The applied mean
compressive traction is \(q_L>0\), so the top compressive resultant is
\(2a q_L\) per unit out-of-plane thickness. Reflection symmetry permits
a quarter-domain calculation for this isotropic problem.

The compatible material inputs, in one consistent nondimensional system,
are \(\phi_{s0}=0.9\), \(K_s=2.5\), mineral shear modulus
\(\mu_s=5/6\), \(K=1\), \(K_f=8\),
\(\bar\rho_{f0}=1\), and \(k/\mu_f=1.5\).
The compliance restriction gives drained shear modulus \(G=0.75\).
The reference Biot coefficient is \(0.6\), and total storage is
\(17/80\). These reference coefficients define the constant-coefficient
analytical problem. The finite-deformation equations instead approach that
problem as \(q_L/K\) tends to zero; they need not reproduce its series
at a finite load.

The undrained post-load state supplies initial displacement and pressure.
At positive time the drained edge has zero pressure. The initial pressure
therefore has an incompatible trace at the drainage boundary, and the
earliest discrete profile contains an unresolved diffusion layer.
We retain the maximum error over all positive output times and also compare
profiles at common fixed positive times under separate spatial and temporal
refinement. Pressure errors are normalized by the initial interior pressure,
and displacement errors by the corresponding final drained displacement.
Neither normalization divides by a transient value approaching zero.

An independent manufactured solution checks off-axis coupling in the
constant reference tangent. On the unit square, prescribe
\begin{align}
 u_1&=U\sin(\pi X_1)\sin(\pi X_2)\sin t,
 \label{eq:fe-mms-u1}\\
 u_2&=U\cos(\pi X_1)\sin(\pi X_2)\sin t,
 \label{eq:fe-mms-u2}\\
 p&=P_0\cos(\pi X_1)\cos(\pi X_2)\sin t,
 \label{eq:fe-mms-pressure}
\end{align}
with \(U=P_0=0.01\). The mineral stiffness in
\eqref{eq:example-mineral-stiffness} is rotated through \(30^\circ\)
about the third material axis. The drained compliance, reference Biot tensor,
and storage are calculated independently from the constitutive restrictions.
Analytic differentiation of the prescribed fields then supplies the body
force and fluid source. Exact displacement and pressure are prescribed on
the boundary, and area-integrated errors are evaluated at a nonzero time.
This test exercises mixed derivatives and off-diagonal pressure coupling;
it does not by itself verify the nonlinear constitutive law.

For a conservative backward-Euler solve, fluid discharge is accumulated
with the same end-step quadrature as the mass residual. The balance error
is normalized by the larger of the mobilized fluid mass and accumulated
absolute discharge, with a stated small floor before drainage begins.
The boundary reaction of the pressure residual provides a discrete
conservative flux. A flux reconstructed from boundary pressure gradients
is evaluated separately because its discretization error is not the
algebraic mass-balance error.

\paragraph{Scope of these results.}
The coupled calculations are an implementation verification of the weak
balances and of the constant reference tangent, and, at finite load and for
prescribed mineral orientations, demonstrations of the constitutive response.
They are not a quantitative verification of the nonlinear finite-deformation
law, and they are not experimental validation: all parameters are synthetic.
The rotated-anisotropy and partial-drainage runs are force-controlled
calculations with a kinematically constrained rigid platen, so their platen
reaction is an equilibrium residual rather than a prescribed displacement.
The partial-drainage runs use a square domain, and their comparison with the
slender Mandel reference geometry is therefore reported as not comparable and
is excluded from the verification claims; only their mass-balance and
platen-equality diagnostics and their maps are reported. The synthetic
parameters and the absence of any laboratory or field comparison are recorded
as a limitation.
